% Springer proceedings working source for CSS 2026.
% Uses Springer's LLNCS proceedings class, available in standard TeX Live.
% If the conference supplies a different Springer Proceedings in Complexity
% class file, the body can be reused with the official class.
\documentclass[runningheads]{llncs}

\usepackage[T1]{fontenc}
\usepackage{amsmath,amssymb,mathtools}
\usepackage{graphicx}
\usepackage{booktabs}
\usepackage{float}
\usepackage{array}
\usepackage{multirow}
\usepackage{url}
\usepackage{xcolor}
\usepackage{tikz}
\usepackage{pgfplots}
\usepackage[linesnumbered,ruled,vlined,noend]{algorithm2e}
\usepackage[hidelinks]{hyperref}
\usetikzlibrary{shapes,calc,positioning,patterns,decorations.markings}
\usepgfplotslibrary{groupplots,statistics,colorbrewer}
\pgfplotsset{compat=1.18}

% Compact equation helpers inherited from the working draft.
\def\EQSP{3pt}
\newcommand{\cgather}[2][\EQSP]{\begingroup\setlength\abovedisplayskip{#1}\setlength\belowdisplayskip{#1}\begin{gather} #2 \end{gather}\endgroup\noindent}
\newcommand{\cgathers}[2][\EQSP]{\begingroup\setlength\abovedisplayskip{#1}\setlength\belowdisplayskip{#1}\begin{gather*} #2 \end{gather*}\endgroup\noindent}
\newcommand{\calign}[2][\EQSP]{\begingroup\setlength\abovedisplayskip{#1}\setlength\belowdisplayskip{#1}\begin{align} #2 \end{align}\endgroup\noindent}
\newcommand{\caligns}[2][\EQSP]{\begingroup\setlength\abovedisplayskip{#1}\setlength\belowdisplayskip{#1}\begin{align*} #2 \end{align*}\endgroup\noindent}
\DeclareMathOperator*{\argmax}{argmax}
\DeclareMathOperator*{\argmin}{arg\,min}

\begin{document}
\raggedbottom
\setlength{\textfloatsep}{12pt plus 2pt minus 2pt}
\setlength{\intextsep}{10pt plus 2pt minus 2pt}
\mainmatter

\title{How Good Is Your Synthetic Social Data? Conditional Alignment and Novelty Across 34 GSS Waves}
\titlerunning{Conditional Alignment and Novelty in Synthetic Social Data}

\author{Hari Dahal\inst{1} \and Ishanu Chattopadhyay\inst{1}\thanks{Corresponding author.}}
\authorrunning{H. Dahal and I. Chattopadhyay}
\institute{University of Kentucky, Lexington, KY, USA\\
\email{ishanu\_ch@uky.edu}}

\maketitle

\begin{abstract}
Synthetic populations and simulated survey respondents are increasingly used in
computational social science, but marginal and task-based validation can miss
the conditional dependencies that organize social attitudes. We introduce a
model-agnostic framework based on full conditional relationships. A
MAP-alignment score compares the probability assigned to each observed response
with the modal conditional probability under the same respondent context; the
complete normalized conditional profile separately induces a metric on strictly
positive finite-state generative processes. We pair MAP alignment with
nearest-real row similarity to measure respondent-level copying. Across 34
General Social Survey waves from 1972--2024, spanning 197--1,273 modeled
variables per wave, a recursive conditional Large Science Model (LSM) achieved
mean MAP alignment $0.8032$, compared with $0.8004$ for the original (control),
for retention $1.0034$ (wave-bootstrap 95\% confidence interval
$[0.9996,1.0079]$). Relative to the independent-marginal null, the excess
MAP-alignment fraction was $1.062$ ($[0.992,1.154]$). Mean nearest-real
similarity was $0.3800$, corresponding to row novelty $0.6200$, whereas the
original (control) had essentially zero novelty. The framework separates conditional alignment from respondent-level
reuse and provides a validation tool for synthetic social data.
\keywords{synthetic populations \and
computational social science \and
survey simulation \and
General Social Survey \and
conditional inference \and
model validation}
\end{abstract}

\section{Introduction and Motivation}
Synthetic social data support data sharing, replication, and population
simulation, but their validity remains difficult to assess. This problem is
increasingly consequential as generative systems are used to construct synthetic
human samples and survey respondents~\cite{argyle2023out,bisbee2024synthetic}.
Records may appear plausible individually while failing to reproduce empirical
response dependencies~\cite{kozlowski2025simulating}. A useful validation
criterion should therefore test population structure without rewarding copies
of observed respondents.

Current practices fall broadly into:
\begin{itemize}
\item \textbf{Likelihood-based metrics} (e.g., ELBO~\cite{kingma2014autoencoding} and perplexity~\cite{jelinek1977perplexity}), requiring the generator to have a tractable joint density;
\item \textbf{Task-based metrics~\cite{esteban2017realvalued}}, such as training a classifier on
      synthetic data and testing on real data, which depend on subjective
      downstream tasks and can obscure structural mismatches;
\item \textbf{Embedding-based metrics} such as FID~\cite{heusel2017gans} or classifier-based
      two-sample tests~\cite{lopezpaz2017revisiting}, which rely on feature extractors unrelated to
      the data domain and  limit interpretability.
\end{itemize}
%% These approaches are not model-agnostic, often depend on architecture
%% choices, and typically do not reveal how the synthetic data preserve
%% (or distort) structural relationships among variables.
Comparisons of componentwise moments or covariance matrices are also
insufficient~\cite{snoke2018general,nowok2016synthpop}: they capture pairwise
linear structure but can miss nonlinear and higher-order dependencies.
Section~\ref{sec:examples} gives two distributions with identical first- and
second-order moments but different conditional structure.

The recursive conditional architecture used by the LSM comparator has supported
digital twins of the infant microbiome~\cite{sizemore2024digital} and influenza
evolution~\cite{wu2026emergenet}. In this benchmark, each wave-specific LSM
models survey variables through coordinate-wise conditional models and generates
complete records by recursive conditional sampling, following the same general
architecture. Here we study evaluation rather than substantive opinion dynamics,
using repeated social surveys as the test bed.

We evaluate conditional structure inferred directly from data. Given real or
synthetic samples, we estimate each coordinate's conditional distribution
$\phi^i(\cdot\mid x^{-i})$, for example using conditional inference trees, and
define the MAP-alignment score
\cgather{
\upsilon(x,i)
=
\frac{\phi^i(x^i \mid x^{-i})}{\max_{y} \phi^i(y\mid x^{-i})},
}
%
For each coordinate, the numerator evaluates the observed value under the
real-data conditional and the denominator is its modal probability in the same
context. Averaging $\upsilon(x,i)$ over records and coordinates yields the
bounded statistic $\Upsilon(D)$. It measures graded agreement with conditional
maximum-a-posteriori predictions~\cite{bishop2006pattern} and requires only
samples and a conditional learner.

Under positivity and compatibility, a complete system of full conditionals
identifies the joint distribution~\cite{brook1964,dobrushin1968,hammersley1971,besag1974}.
This motivates a one-sided MAP-alignment comparison and a symmetric distance
between full conditional profiles. Unlike aggregate utility measures
\cite{woo2009global,elemam2020seven,elemam2022utility,dankar2022multidimensional,kaabachi2025scoping},
we assess conditional structure and proximity to individual real records
separately.

\section{Formal Setup for Conditional Analysis}
Let
\cgather{
X=(X^1,\dots,X^N)
}
take values in the finite product space
\cgather{
\mathcal{X}=\mathcal{X}^1\times\cdots\times\mathcal{X}^N.
}
Let $P$ be a strictly positive data-generating distribution. For coordinate
$i$, define the full conditional
\cgather{
P_i(x^i\mid x^{-i})
:=P(X^i=x^i\mid X^{-i}=x^{-i}).
}
A learned conditional kernel is denoted by
$\phi^i(\cdot\mid x^{-i})$. In practice, a small $\varepsilon$-floor may be
added and the conditional vector renormalized. Positivity is used together
with compatibility; it does not by itself guarantee that an arbitrary family
of kernels corresponds to a joint distribution.

\section{MAP-Alignment Functional}
For any model $\{\phi^i\}$ and sample $x\in\mathcal{X}$, define
\cgather{
    \upsilon(x,i)
    :=
    \frac{\phi^i(x^i\mid x^{-i})}
         {\max_{y\in \mathcal{X}^i}
          \phi^i(y\mid x^{-i})},
}
which equals $1$ exactly when $x^i$ is a maximizer of the model conditional.

Given a dataset $D=\{x_{k}\}_{k=1}^M$, define
\cgather{
    \Upsilon(D)
    :=
    \frac{1}{MN}
    \sum_{k=1}^M\sum_{i=1}^N
    \upsilon(x_{k},i),
}
an empirical estimate of
\cgather{
    \Upsilon_\phi(P)
    =
    \mathbb{E}_{X\sim P}
    \left[
        \frac{1}{N}\sum_{i=1}^N \upsilon(X,i)
    \right].
}
\subsection{Behavior Under Exact Conditionals}

\begin{lemma}[MAP-Alignment Under Exact Conditionals]
\label{lem:map}
Assume $\phi^i = P_i$ for all $i$.
Fix $i$ and $x^{-i}$, and let $p_j := P_i(j\mid x^{-i})$,
$p_{\max} := \max_j p_j$.
If $X^i \sim P_i(\cdot\mid x^{-i})$, then
\cgather{
    \upsilon(X,i) = \frac{P_i(X^i\mid x^{-i})}{p_{\max}},
\qquad
    \mathbb{E}[\upsilon(X,i)\mid x^{-i}]
    =
    \frac{1}{p_{\max}}\sum_j p_j^2.
}
Moreover:
\begin{itemize}
\item $\upsilon(X,i)=1$ iff $X^i\in\arg\max_j p_j$.
\item If $P_i(\cdot\mid x^{-i})$ is uniform on its support,
      then $\mathbb{E}[\upsilon(X,i)\mid x^{-i}] = 1$.
\item If some $p_j < p_{\max}$, then
      $\mathbb{E}[\upsilon(X,i)\mid x^{-i}] < 1$.
\end{itemize}
\end{lemma}

\begin{proof}
Immediate from the definition and the fact that
$\sum_j p_j^2 \le p_{\max}\sum_j p_j = p_{\max}$.
\end{proof}

Thus, under exact conditionals, $\upsilon(X,i)$ is a normalized
conditional-likelihood score. Its population target is
\cgather{
    \Upsilon_{\mathrm{oracle}}(P)
    :=
    \mathbb{E}_{X\sim P}
    \left[
        \frac{1}{N}\sum_{i=1}^N
        \frac{P_i(X^i\mid X^{-i})}
             {\max_y P_i(y\mid X^{-i})}
    \right].
    \label{eq:oracle_ratio}}
This target is generally smaller than one when the exact conditional is
nonuniform. Consequently, synthetic MAP alignment is assessed by agreement with a
real-data reference score, not by requiring $\Upsilon\to1$.

\section{Brook--Dobrushin Factorization and Identifiability}
Brook's lemma~\cite{brook1964} and related consistency results~\cite{dobrushin1968}
show how a compatible system of full conditionals identifies a strictly
positive joint distribution.

Assume
\cgather{
    \phi^i(x^i\mid x^{-i}) = P_i(x^i\mid x^{-i})
    \quad \text{whenever } P(x)>0.
}
Choose a reference configuration $x^\circ$ with $P(x^\circ)>0$.
Define the interpolating sequence
\cgather{
    x^{(0)} = x^\circ, \qquad
    x^{(i)} =
    (x^1,\dots,x^i, x^{i+1,\circ},\dots,x^{N,\circ}),
}
so $x^{(N)} = x$.
By repeated conditioning,
\cgather{
    \frac{P(x^{(i)})}{P(x^{(i-1)})}
    =
    \frac{P_i(x^i\mid x^{<i}, x^{>i,\circ})}
         {P_i(x^{i,\circ}\mid x^{<i}, x^{>i,\circ})}.
}
Multiplying yields the Brook factorization~\cite{brook1964}:
\cgather{
    \frac{P(x)}{P(x^\circ)}
    =
    \prod_{i=1}^N
    \frac{P_i(x^i\mid x^{<i}, x^{>i,\circ})}
         {P_i(x^{i,\circ}\mid x^{<i}, x^{>i,\circ})}.
}
Replacing $P_i$ by $\phi^i$ on the support gives an explicit
reconstruction of $P$ from the conditionals.

\begin{theorem}[Identification from compatible full conditionals]
\label{thm:brook}
Let $P$ and $Q$ be strictly positive distributions on the same finite product
space. If their full conditionals agree for every coordinate and every
configuration, then $P=Q$.
\end{theorem}

\begin{proof}
Brook's factorization reconstructs each joint probability ratio relative to a
fixed reference configuration from the full conditional system. Equality of
all full conditionals therefore gives identical probability ratios under $P$
and $Q$; normalization gives $P=Q$~\cite{brook1964,dobrushin1968}. The
compatibility assumption is essential: positivity alone does not guarantee that
an arbitrary collection of conditional kernels corresponds to a joint law.
\end{proof}

\section{Evaluating Generators via \texorpdfstring{$\Upsilon$}{Upsilon}}
For any model $\{\phi^i\}$,
\cgather{
    \Upsilon_\phi(P)
    =
    \mathbb{E}_{X\sim P}
    \left[
        \frac{1}{N}\sum_{i=1}^N
        \frac{\phi^i(X^i\mid X^{-i})}
             {\max_y \phi^i(y\mid X^{-i})}
    \right].
}
Let $D_{\mathrm{test}}$ be i.i.d.\ data from $P$. If $\phi^i_n\to P_i$
pointwise on the support and satisfy strict positivity, then dominated
convergence yields
\cgather{
    \Upsilon_{\phi_n}(P)
    \to
    \Upsilon_{\mathrm{oracle}}(P).
}
\begin{corollary}[Conditional-Kernel Convergence and MAP-Alignment Consistency]
Let $\phi^i_n$ be strictly positive kernels converging to $P_i$ on
$\mathrm{supp}(P)$. Let $\widetilde P_n$ denote the joint compatible
with $\{\phi^i_n\}$. Then:
\begin{itemize}
\item $\widetilde P_n \to P$ pointwise (equivalently, in total variation on the finite state space).
\item For any sequence of test sets $D_{\mathrm{test}}$ with
      $|D_{\mathrm{test}}|\to\infty$, the empirical scores
      $\Upsilon(D_{\mathrm{test}};\phi_n)$ converge in probability to
      $\Upsilon_{\mathrm{oracle}}(P)$.
\end{itemize}
Both conclusions follow from the assumed convergence of the full conditional
kernels. Convergence, or numerical agreement, of the scalar MAP-alignment score
alone does not imply convergence of the conditional kernels or recovery of the
joint law. Systematically low MAP alignment can nevertheless reveal disagreement
with the evaluator's conditional modes even when marginal summaries look
satisfactory.
\end{corollary}

\section{Conditional Inference, MAP-Alignment Uncertainty, and a Metric on Underlying Processes}

\subsection{Inference of Conditionals Using Conditional Learners}

Let $D=\{x^k\}_{k=1}^n$ be a training dataset in a finite product space. We
construct a family of full conditional models
$\Phi^{(n)}=\{\varphi^{i,n}(\cdot\mid x^{-i})\}_{i=1}^N$ by solving one
supervised prediction problem for each coordinate. In the experiments, these
models are implemented using conditional inference trees, whose unbiased
permutation-based split selection is described in
\cite{hothorn2006,hothorn2006party,strobl2007}.

The theory does not require a tree-specific convergence exponent. We instead
state the estimation assumption explicitly: for each coordinate,
\begin{equation}
\begin{aligned}
&\sup_{x^{-i},x^i}
\left|\varphi^{i,n}(x^i\mid x^{-i})-P_i(x^i\mid x^{-i})\right|
\leq \delta_{i,n},\\
&\delta_n:=\max_i\delta_{i,n}\xrightarrow{p}0.
\end{aligned}
\label{eq:conditional_error}
\end{equation}
A small $\varepsilon$-floor may be applied before renormalization to maintain
strict positivity on the effective support. When the learned kernels are
compatible, the resulting full conditional system identifies a unique joint
law by Theorem~\ref{thm:brook}.

\subsection{MAP-Alignment and Its Finite-Sample Uncertainty}

Given a trained conditional family $\Phi^{(n)}$, define
\cgather{
\upsilon_{\Phi^{(n)}}(x,i)
=
\frac{\varphi^{i,n}(x^i\mid x^{-i})}
{\max_{y\in\mathcal{X}^i}\varphi^{i,n}(y\mid x^{-i})}
\in[0,1].
\label{eq:upsilon_definition}
}
For an independent test dataset $D_{\mathrm{test}}=\{x_k\}_{k=1}^M$, let
\cgather{
Z_k:=\frac{1}{N}\sum_{i=1}^N
\upsilon_{\Phi^{(n)}}(x_k,i),
\qquad
\hat\Upsilon_{\Phi^{(n)}}(D_{\mathrm{test}})
:=\frac{1}{M}\sum_{k=1}^M Z_k.
\label{eq:upsilon_empirical}
}
The coordinates within a record may be dependent, but the row scores $Z_k$ are
independent and bounded in $[0,1]$ under i.i.d.\ test sampling. Hence
Hoeffding's inequality~\cite{hoeffding1963} gives
\cgather{
\Pr\left(
\left|\hat\Upsilon_{\Phi^{(n)}}-\mathbb{E}Z_1\right|>\epsilon
\right)
\leq 2\exp(-2M\epsilon^2).
\label{eq:upsilon_hoeffding}
}
Because each finite conditional vector has modal probability at least
$1/|\mathcal{X}^i|$, the normalization in
Eq.~\eqref{eq:upsilon_definition} is Lipschitz in the conditional vector. Combining test-sample variation with
Eq.~\eqref{eq:conditional_error} yields
\cgather{
\hat\Upsilon_{\Phi^{(n)}}(D_{\mathrm{test}})
=
\Upsilon_{\mathrm{oracle}}(P)
+O_p\!\left(M^{-1/2}+\delta_n\right).
\label{eq:upsilon_full_uncertainty}
}

\subsection{A Metric on Underlying Generative Processes}

For a strictly positive process $P$ with full conditionals $P_i$, define its population $\upsilon$-profile:
\cgather{
u_P(x,i)
=
\frac{
P_i(x^i\mid x^{-i})
}{
\max_{y\in\mathcal{X}^i} P_i(y\mid x^{-i})
}.
}
Fix a reference probability measure $\mu$ on $\mathcal{X}$ with full
support, and sample the coordinate index uniformly from
$\{1,\dots,N\}$. Define the distance
\cgather{
d_{\mu}(P_1,P_2)
=
\mathbb{E}_{X\sim\mu}
\left[
\frac{1}{N}\sum_{i=1}^{N}
\lvert u_{P_1}(X,i) - u_{P_2}(X,i)\rvert
\right].
\label{eq:true_metric}
}

\begin{theorem}
Let $\mu$ have full support on the finite product space $\mathcal{X}$.
Then $d_{\mu}$ is a metric on the set of strictly positive probability
distributions on $\mathcal{X}$.
\end{theorem}

\begin{proof}
Nonnegativity and symmetry are immediate. The triangle inequality follows from the scalar inequality
$\lvert a-c\rvert \le \lvert a-b\rvert + \lvert b-c\rvert$ applied pointwise and averaged over coordinates and integrated with respect to $\mu$.
For identity of indiscernibles, if $d_{\mu}(P_1,P_2)=0$, then
$u_{P_1}(x,i)=u_{P_2}(x,i)$ for every $(x,i)$ because $\mathcal{X}$ is
finite and $\mu$ has full support. For each fixed
$(x^{-i},i)$, the conditional distribution is recovered from its normalized
profile by
$P_i(x^i\mid x^{-i})=u_P(x,i)/\sum_{z\in\mathcal{X}^i}u_P((z,x^{-i}),i)$.
Thus all full conditionals agree. Theorem~\ref{thm:brook} then implies $P_1=P_2$.
\end{proof}

Given empirical datasets $D_1,D_2$, train conditional generator families
$G_1,G_2$ yielding conditional families $\Phi^{(1)}$ and $\Phi^{(2)}$.
Independently draw a common evaluation sample
$E_{\mu}=\{x^{(k)}\}_{k=1}^{M}$ i.i.d.\ from the fixed reference measure
$\mu$. Define the empirical distance
\cgather{
\widehat d_{\mu}(D_1,D_2;E_{\mu})
=
\frac{1}{MN}
\sum_{k=1}^{M}
\sum_{i=1}^{N}
\left\lvert
u_{G_1}(x^{(k)},i)-u_{G_2}(x^{(k)},i)
\right\rvert.
\label{eq:empirical_metric}
}

For notational convenience, let $d_{\mu}(G_1,G_2)$ denote
Eq.~\eqref{eq:true_metric} evaluated using the two learned normalized profiles.

\begin{theorem}
Assume the conditional estimators satisfy \eqref{eq:conditional_error} with
uniform errors $\epsilon_1,\epsilon_2$. Then
\cgather{
\big|
\widehat d_{\mu}(D_1,D_2;E_{\mu})-d_{\mu}(P_1,P_2)
\big|
=
O_p\!\left(
M^{-1/2}
+
\epsilon_1
+
\epsilon_2
\right).
}
\end{theorem}

\begin{proof}
By the triangle inequality,
\calign{
\left\lvert \widehat d_{\mu}-d_{\mu}(P_1,P_2)\right\rvert
&\leq
\left\lvert \widehat d_{\mu}-d_{\mu}(G_1,G_2)\right\rvert\\
&\quad+
\left\lvert d_{\mu}(G_1,G_2)-d_{\mu}(P_1,P_2)\right\rvert.
}
For the first term, the row-level quantities
\cgather{
Z_k=\frac{1}{N}\sum_{i=1}^{N}
\left\lvert u_{G_1}(x^{(k)},i)-u_{G_2}(x^{(k)},i)\right\rvert
}
are i.i.d.\ and bounded in $[0,1]$ because $E_{\mu}$ is independent of the
training datasets. Hoeffding's inequality therefore gives the
$O_p(M^{-1/2})$ term. The second term is controlled by the Lipschitz
continuity of $u(\varphi)=\varphi/\max_y\varphi(y)$ on finite
probability simplices, together with \eqref{eq:conditional_error}.
\end{proof}
\subsection{One-Sided MAP-Alignment Comparison Against Held-Out Real Data}

Let $D_{\mathrm{fit}}$ be an evaluator-training dataset independent of both
the real reference dataset $D_{\mathrm{ref}}$ and the synthetic dataset
$D_{\mathrm{syn}}$. Fit a conditional family $\Phi^{(\mathrm{real})}$ on
$D_{\mathrm{fit}}$, and evaluate both the real reference records and synthetic
records under that same fixed conditional family:
\cgather{
\Upsilon_{\mathrm{ref}}
=\hat\Upsilon_{\Phi^{(\mathrm{real})}}(D_{\mathrm{ref}}),
\qquad
\Upsilon_{\mathrm{syn}}
=\hat\Upsilon_{\Phi^{(\mathrm{real})}}(D_{\mathrm{syn}}).
}
We report the signed MAP-alignment gap and retention ratio
\cgather{
\Delta\Upsilon:=\Upsilon_{\mathrm{syn}}-\Upsilon_{\mathrm{ref}},
\qquad
R_{\Upsilon}:=\frac{\Upsilon_{\mathrm{syn}}}{\Upsilon_{\mathrm{ref}}}.
\label{eq:fidelity_gap_retention}
}
A generator matching the reference in average MAP alignment should have
$\Delta\Upsilon$ near zero, or equivalently $R_{\Upsilon}$ near one. Negative
gaps indicate lower average MAP alignment;
positive gaps do not by themselves establish superiority because finite-sample
estimation and mode concentration can increase the normalized score. When an
independent evaluator-training dataset is unavailable, cross-fitting may instead
be used to avoid evaluating real records under a model trained on those same
records.

\subsection{Consistency of the Dataset Distance as a Metric on Processes}
\begin{theorem}[Convergence to the Population Metric]
Let $D_1$ and $D_2$ be drawn i.i.d.\ from strictly positive processes
$P_1$ and $P_2$, and let $E_{\mu}$ be an independent i.i.d.\ sample from the
fixed full-support reference measure $\mu$. Suppose the learned conditional kernels $\widehat P_{k,i}$ satisfy
\calign{
&\max_{1\leq i\leq N}\sup_{x^{-i},x^i}
\left\lvert
\widehat P_{k,i}(x^i\mid x^{-i})-P_{k,i}(x^i\mid x^{-i})
\right\rvert\\
&\hspace{5em}\leq \epsilon_k,
\qquad k\in\{1,2\}.
}
with $\epsilon_k\to 0$ in probability as $|D_k|\to\infty$, and suppose
$M=|E_{\mu}|\to\infty$. Then
\cgather{
\widehat d_{\mu}(D_1,D_2;E_{\mu})
\xrightarrow{p}
d_{\mu}(P_1,P_2).
\label{eq:limit_metric}
}
\end{theorem}

\begin{proof}
From the previous theorem we have
\cgather{
\big|
\widehat d_{\mu}(D_1,D_2;E_{\mu})-d_{\mu}(P_1,P_2)
\big|
=
O_p\!\left(
M^{-1/2}
+
\epsilon_1
+
\epsilon_2
\right).
}
As $M\to\infty$ and $|D_k|\to\infty$, both $M^{-1/2}$ and $\epsilon_k$ converge to zero in probability, which implies \eqref{eq:limit_metric}.
\end{proof}

\begin{corollary}
If $P_1=P_2$, then
\cgather{
\widehat d_{\mu}(D_1,D_2;E_{\mu})\xrightarrow{p}0.
}
If $P_1\neq P_2$, then
\cgather{
\widehat d_{\mu}(D_1,D_2;E_{\mu})
\xrightarrow{p}d_{\mu}(P_1,P_2)>0.
}
\end{corollary}

\subsection{Cross-Fitted Evaluation When Independent Training Data Are Unavailable}

Algorithm~\ref{alg:crossfit} gives a general cross-fitted implementation for
settings in which a separate evaluator-training dataset is unavailable. In the
GSS benchmark reported here, the wave-specific conditional evaluators were
trained on data disjoint from the corresponding real reference records and from all generated samples. The reported GSS scores are
therefore out-of-sample with respect to evaluator fitting and do not require the
cross-fitting procedure in Algorithm~\ref{alg:crossfit}.

\begin{algorithm}[t]
\caption{Cross-fitted conditional MAP alignment and row novelty}
\label{alg:crossfit}
\DontPrintSemicolon
\SetKwInOut{Input}{Input}
\SetKwInOut{Output}{Output}
\Input{Real data $D_{\mathrm{real}}$; synthetic data $D_{\mathrm{syn}}$;
$K$ folds; conditional learner $\mathsf{Learn}$; level $\alpha$.}
\Output{$\hat\Upsilon_{\mathrm{ref}}$, $\hat\Upsilon_{\mathrm{syn}}$,
$\Delta\hat\Upsilon$, $R_{\Upsilon}$, and row novelty.}
Partition $D_{\mathrm{real}}$ into folds $F_1,\dots,F_K$\;
\For{$k\leftarrow1$ \KwTo $K$}{
Fit $\Phi_k\leftarrow\mathsf{Learn}(D_{\mathrm{real}}\setminus F_k)$\;
Score each $x\in F_k$ using $\Phi_k$ and store its row-average MAP alignment\;
Score an assigned subset of $D_{\mathrm{syn}}$ using the same $\Phi_k$\;
}
Average the stored real and synthetic row scores to obtain
$\hat\Upsilon_{\mathrm{ref}}$ and $\hat\Upsilon_{\mathrm{syn}}$\;
Set $\Delta\hat\Upsilon\leftarrow
\hat\Upsilon_{\mathrm{syn}}-\hat\Upsilon_{\mathrm{ref}}$ and
$R_{\Upsilon}\leftarrow
\hat\Upsilon_{\mathrm{syn}}/\hat\Upsilon_{\mathrm{ref}}$\;
For each synthetic row, compute nearest-real similarity $\eta_j$ and report
row novelty $1-m^{-1}\sum_j\eta_j$\;
Use a paired bootstrap over independent datasets or survey waves for the final
confidence interval\;
\Return all summary statistics\;
\end{algorithm}

For a fixed fitted conditional family and $M$ independent evaluated rows, the
row-level Hoeffding radius at confidence $1-\delta$ is
\begin{equation}
r_M(\delta)=\sqrt{\frac{1}{2M}\log\frac{2}{\delta}}.
\end{equation}
At $\delta=0.05$, the radii are $0.1358$ for $M=100$ and $0.0429$ for
$M=1000$. Conditional-estimation uncertainty must be handled separately through
data splitting, repeated fitting, or a higher-level bootstrap.

\section{Low-Order Moment Matching is Insufficient: Two Illustrative Examples}\label{sec:examples}

Two examples in $\mathbb{R}^3$ show that identical first- and second-order
moments do not imply similar conditional structure. We use the
continuous-density analogue of the finite-state score, replacing the maximum by
an essential supremum. A formal extension to continuous and mixed-type data is
left for future work.

\subsection{Example 1: Uniform vs.\ Gaussian with Identical Moments}

Consider the following two distributions on $\mathbb{R}^3$:
\cgather{
X^{(U)} \sim \mathrm{Uniform}([-\sqrt{3},\sqrt{3}]^3),
\qquad
X^{(G)} \sim \mathcal{N}_3(0,I_3).
}

Both satisfy
\cgather{
\mathbb{E}[X^{(U)}]
=
\mathbb{E}[X^{(G)}]
=
0,
\qquad
\mathrm{Cov}(X^{(U)})
=
\mathrm{Cov}(X^{(G)})
=
I_3.
}
Hence all marginal means, variances, and the full covariance matrix
coincide.

Yet their conditional structures differ sharply. For $X^{(U)}$, each
coordinate has a flat conditional density on $[-\sqrt{3},\sqrt{3}]$,
yielding
\cgather{
\upsilon(X^{(U)}, i) = 1
\quad\text{for almost every sample}.
}
For $X^{(G)}$, the $i$th coordinate conditional is
$X^{(G)}_i \mid X^{(G)}_{-i} \sim \mathcal{N}(0,1)$, giving
\cgather{
\upsilon(X^{(G)}, i)
=
\exp\!\left(-\tfrac12 X^{(G)}_i{}^2\right),
}
with mean exactly $1/\sqrt{2}\approx0.7071$. Thus, although the datasets match in mean
and covariance, the MAP-alignment values differ substantially:
\cgather{
\Upsilon(X^{(U)}) = 1,
\qquad
\Upsilon(X^{(G)}) = \frac{1}{\sqrt{2}}.
}
Low-order moments fail to detect this discrepancy.

\subsection{Example 2: Matching Non-Identity Covariances with Distinct Conditional Structure}

To demonstrate that the limitation persists even when the covariance
matrix is nontrivial, apply the same invertible linear map
\cgather{
L
=
\begin{pmatrix}
1 & \rho & 0\\
0 & 1     & \rho\\
0 & 0     & 1
\end{pmatrix},
\qquad 0 < \rho < 1,
}
to both datasets and define
\cgather{
Y^{(U)} = L X^{(U)},
\qquad
Y^{(G)} = L X^{(G)}.
}

Both transformed datasets have the \emph{same} non-identity covariance
matrix
\cgather{
\mathrm{Cov}(Y^{(U)})
=
\mathrm{Cov}(Y^{(G)})
=
\Sigma
=
L L^{\!\top},
}
and share identical columnwise means and variances.

Their conditional distributions nevertheless differ. Since $Y^{(U)}$ is the
image of a uniform cube under a linear shear, each
conditional $Y^{i,(U)} \mid Y^{-i,(U)}$ is uniform on a finite
interval (given by the intersection of a line with a parallelepiped),
implying
\cgather{
\upsilon(Y^{(U)}, i) = 1
\quad\text{for almost every sample}.
}

Conversely, $Y^{(G)} \sim \mathcal{N}_3(0,\Sigma)$ has
linear--Gaussian conditionals. If
$Y^{i,(G)} \mid Y^{-i,(G)} \sim \mathcal{N}(m_i(y^{-i}), \sigma_i^2)$,
then
\cgather{
\upsilon(Y^{(G)}, i)
=
\exp\!\left(
-\tfrac12 (y^i - m_i(y^{-i}))^2 / \sigma_i^2
\right),
}
and the standardized conditional residual is $\mathcal{N}(0,1)$, so its mean is exactly $1/\sqrt{2}\approx0.7071$.

Thus, even though $Y^{(U)}$ and $Y^{(G)}$ agree in all first- and
second-order statistics, their conditional behavior differs markedly,
and the MAP-alignment statistic again reveals the discrepancy:
\cgather{
\Upsilon(Y^{(U)}) = 1,
\qquad
\Upsilon(Y^{(G)}) = \frac{1}{\sqrt{2}}.
}

\subsection{Implications}

Thus, matching means and covariance does not determine conditional structure.
MAP alignment detects the discrepancy in both examples.


\section{Applications: Comparing Synthesizers}\label{sec:applications}

\subsection{Row-Level Memorization Diagnostic}

MAP alignment does not detect row reuse: a dataset can score highly while
copying observed records. We therefore add a row-level novelty diagnostic.

Let $D_{\mathrm{real}}=\{x^{(1)},\ldots,x^{(n)}\}$ and
$D_{\mathrm{syn}}=\{y^{(1)},\ldots,y^{(m)}\}$ be real and synthetic datasets
defined over the same categorical (or discretized) feature space. Using a
concatenated one-hot embedding $f(\cdot)$, we define cosine similarity
\cgather{
s(u,v)
\;=\;
\frac{\langle f(u), f(v)\rangle}
{\|f(u)\|_2\,\|f(v)\|_2},
} which, for categorical data, reduces to the fraction of coordinates on which two
records agree.

For each synthetic record $y^{(j)}$, we define its nearest-real similarity
\cgather{
\eta_j
\;=\;
\max_{1 \le i \le n}
s\!\left(x^{(i)}, y^{(j)}\right).
}
We summarize row-level reuse via the \emph{lack-of-novelty score}
\cgather{
\mathcal{N}(D_{\mathrm{syn}})
\;=\;
\frac{1}{m}\sum_{j=1}^m \eta_j,
}
with larger values indicating lower novelty. In particular,
$\mathcal{N}=1$ corresponds to exact row reuse (e.g., bootstrap resampling),
while smaller values indicate increasing deviation from any individual real
record. Tail statistics of $\{\eta_j\}$ (e.g., the 95th percentile)
capture worst-case near-copying behavior.

The diagnostics are complementary: $\Upsilon(D)$ measures conditional
alignment, whereas $\mathcal{N}(D_{\mathrm{syn}})$ measures proximity to
specific real records. Nearest-row novelty is a memorization diagnostic, not a
formal privacy guarantee; disclosure risk requires separate membership, linkage,
or attribute-inference analysis. The target is $\Delta\Upsilon\to0$
(equivalently $R_{\Upsilon}\to1$) with $\mathcal{N}$ well below one.

\subsection{Practical Use of the Diagnostics}

The diagnostics answer different questions and should be selected accordingly.
MAP alignment is appropriate when the primary concern is whether synthetic
records preserve the conditional response structure represented by a held-out
real-data evaluator. In that setting, practitioners should report the signed
MAP-alignment gap, the raw reference score, and, when an independent-marginal
null is available, the zero-anchored excess-alignment fraction $\Lambda$.

Nearest-real similarity is appropriate when the concern is copying or near reuse
of individual records. It should be reported together with exact duplicate rates
and upper-tail summaries, such as the 95th or 99th percentile, rather than only
a mean. It is not a substitute for formal disclosure-risk analysis.

Neither diagnostic replaces marginal checks or application-specific validation.
When the synthetic data will support a particular estimand, classifier, or policy
simulation, MAP alignment and novelty should be combined with relevant marginal,
task-based, and, where necessary, privacy metrics. A single scalar score is not
sufficient for all uses of synthetic social data.


\section{Empirical Evaluation on Repeated GSS Cross-Sections}
\label{sec:gss_results}

\subsection{Benchmark Design}

We analyzed 34 General Social Survey (GSS) waves spanning
1972--2024~\cite{gss2024}. Each wave was split before model fitting: 50\% of
respondents were used to train the wave-specific LSM and conditional evaluator,
and the remaining 50\% were held out from LSM fitting for comparator
construction, the real-data control, and all reported evaluation. Missing and
nonresponse entries were preserved as missing throughout LSM fitting and
evaluation rather than recoded as substantive response levels or statistically
imputed. For nearest-row similarity, missing entries were represented by a
common missing marker solely for one-hot computation. Across waves, the data
contained 71,667 respondent records and 197--1,273 modeled variables per wave
(Table~\ref{tab:gss_data}).

\begin{table}[t]
\centering
\caption{GSS data characteristics and analysis split.}
\label{tab:gss_data}
\small
\setlength{\tabcolsep}{3.5pt}
\begin{tabular}{@{}lr@{}}
\toprule
Characteristic & Value \\
\midrule
Analyzed waves & 34 \\
Survey years covered & 1972--2024 \\
Respondent records & 71,667 \\
Respondents per wave & 1,372--4,510 \\
LSM/evaluator training split & 50\% per wave \\
Held-out analysis split & 50\% per wave \\
Modeled variables per wave & 197--1,273 \\
Evaluated records per condition & 1,000 \\
\bottomrule
\end{tabular}
\end{table}

We compared LSM, a restricted Chow--Liu hybrid~\cite{chow1968},
CTGAN~\cite{xu2019ctgan}, an independent marginal baseline, and the original
(control). Each condition produced 1,000 records and
was scored by the same fixed, wave-specific evaluator. Waves were the paired
comparison units. We report the signed gap
$\Upsilon_{g,w}-\Upsilon_{\mathrm{control},w}$, the retention ratio
$\Upsilon_{g,w}/\Upsilon_{\mathrm{control},w}$, and row novelty. Confidence
intervals are nonparametric bootstraps over the 34 waves. Because GSS waves are
temporally ordered repeated cross-sections rather than exchangeable draws from a
population of survey years, these intervals summarize stability across the
observed waves and are not population-level temporal confidence intervals.
Numerical similarity values were clipped to $[0,1]$.

Because $\Upsilon$ is bounded but not zero-anchored, we additionally report the
fraction of alignment recovered above the independent-marginal null. Let
$\bar{\Upsilon}_g$ denote the mean MAP-alignment score of generator $g$ across
waves. We define
\begin{equation}
\Lambda_g
=
\frac{\bar{\Upsilon}_g-\bar{\Upsilon}_{\mathrm{ind}}}
{\bar{\Upsilon}_{\mathrm{control}}-\bar{\Upsilon}_{\mathrm{ind}}}.
\label{eq:excess_alignment}
\end{equation}
Thus, the independent-marginal baseline is anchored at $\Lambda=0$ and the
original (control) at $\Lambda=1$. Bootstrap intervals resample waves jointly
and recompute the complete ratio within each replicate. Table~\ref{tab:gss_benchmark}
summarizes the cross-wave results.

\begin{table}[t]
\centering
\caption{Cross-wave benchmark summary over 34 GSS waves. The signed gap is
$\Upsilon_{\mathrm{generator}}-\Upsilon_{\mathrm{control}}$.
The excess-alignment fraction $\Lambda$ is defined in
Eq.~\eqref{eq:excess_alignment}. Confidence intervals are nonparametric 95\%
bootstrap intervals over waves; the null and control entries in the $\Lambda$
column are fixed anchors. Row novelty is one minus mean nearest-real
similarity.}
\label{tab:gss_benchmark}
\small
\setlength{\tabcolsep}{3.5pt}
\resizebox{\textwidth}{!}{%
\begin{tabular}{lccccc}
\toprule
Generator & Mean $\Upsilon$ & Signed gap (95\% CI) &
Excess $\Lambda$ (95\% CI) & Nearest-real similarity & Row novelty \\
\midrule
LSM & 0.8032 & 0.0027 $[-0.0004,0.0062]$ &
1.062 $[0.992,1.154]$ & 0.3800 & 0.6200 \\
Chow--Liu hybrid & 0.7872 & $-0.0132$ $[-0.0159,-0.0106]$ &
0.698 $[0.661,0.736]$ & 0.6560 & 0.3440 \\
CTGAN & 0.7630 & $-0.0375$ $[-0.0415,-0.0332]$ &
0.143 $[0.068,0.216]$ & 0.5665 & 0.4335 \\
Independent baseline & 0.7567 & $-0.0437$ $[-0.0483,-0.0390]$ &
0.000 (anchor) & 0.5619 & 0.4381 \\
Original (control) & 0.8004 & 0.0000 & 1.000 (anchor) & 1.0000 & 0.0000 \\
\bottomrule
\end{tabular}%
}
\end{table}

\subsection{MAP Alignment}

Across the 34 waves, LSM attained mean $\Upsilon=0.8032$ versus $0.8004$
for the original (control). The signed difference was $0.0027$
(wave-bootstrap 95\% CI $[-0.0004,0.0062]$), and retention was $1.0034$ (95\%
CI $[0.9996,1.0079]$). LSM exceeded the control in 22 waves, but the interval
includes parity; we therefore interpret this as control-level MAP alignment,
not superiority or equality of the full conditional distributions.

The restricted Chow--Liu hybrid retained $0.9834$ of control alignment
(95\% CI $[0.9801,0.9866]$), CTGAN retained $0.9530$ (95\% CI
$[0.9479,0.9584]$), and the independent baseline retained $0.9451$ (95\% CI
$[0.9394,0.9510]$). These ratios are not zero-anchored. On the excess-alignment
scale in Eq.~\eqref{eq:excess_alignment}, LSM attained
$\Lambda_{\mathrm{LSM}}=1.062$ (95\% CI $[0.992,1.154]$), the restricted
Chow--Liu hybrid attained $0.698$ ($[0.661,0.736]$), and CTGAN attained $0.143$
($[0.068,0.216]$). The LSM interval includes one, so its point estimate above
one indicates consistency with control-level alignment, not superiority to real
data. Figure~\ref{fig:gss_benchmark_four_panel}(b) shows the paired wave-level
signed-gap distributions, while Fig.~\ref{fig:gss_benchmark_four_panel}(d)
shows MAP-alignment retention across GSS waves.

\subsection{MAP Alignment Without Row Reuse}

The original (control) had nearest-real similarity essentially one. Mean
similarity was $0.3800$ for LSM, $0.6560$ for the restricted Chow--Liu hybrid,
$0.5665$ for CTGAN, and $0.5619$ for the independent baseline, corresponding to
novelty $0.6200$, $0.3440$, $0.4335$, and $0.4381$, respectively.
Figure~\ref{fig:gss_benchmark_four_panel}(c) shows the wave-level novelty
distributions; Fig.~\ref{fig:gss_benchmark_four_panel}(a) combines mean novelty
with excess MAP alignment above the independent-marginal null. LSM therefore
occupies the favorable region of the archived alignment--novelty plane.


\begin{figure}[tbp]
  \centering
  % The boxplot command below exposes all visible components directly.
% PGFPlots computes quartiles and 1.5-IQR whiskers from the raw CSV column;
% no mean marker is drawn.
\newcommand{\ControlledBoxPlot}[4]{%
  \addplot[
    boxplot/draw direction=y,
    boxplot={
      draw position=#1,
      box extend=0.37,
      whisker range=2.5
    },
    draw=#2,
    fill=#2!9,
    line width=0.55pt,
    boxplot/every box/.style={draw=#2, fill=#2!9, line width=0.55pt},
    boxplot/every whisker/.style={draw=#2, line width=0.50pt},
    boxplot/every median/.style={draw=black, line width=0.85pt},
    mark=o,
    mark size=1.2pt,
    mark options={draw=#2, fill=white, line width=0.35pt}
  ] table[y=#4, col sep=comma] {#3};%
}
% Data files.
\newcommand{\WaveLSM}{data/wave_lsm.csv}
\newcommand{\WaveChow}{data/wave_chow_liu.csv}
\newcommand{\WaveCTGAN}{data/wave_ctgan.csv}
\newcommand{\WaveBase}{data/wave_baseline.csv}
\newcommand{\WaveCtrl}{data/wave_original_sample_control.csv}
\newcommand{\SumLSM}{data/summary_lsm.csv}
\newcommand{\SumChow}{data/summary_chow_liu.csv}
\newcommand{\SumCTGAN}{data/summary_ctgan.csv}
\newcommand{\SumBase}{data/summary_baseline.csv}
\newcommand{\SumCtrl}{data/summary_original_sample_control.csv}

% Consistent generator colors.
\definecolor{lsmcol}{RGB}{31,119,180}
\definecolor{chowcol}{RGB}{230,126,34}
\definecolor{ctgancol}{RGB}{46,139,87}
\definecolor{basecol}{RGB}{192,57,43}
\definecolor{ctrlcol}{RGB}{117,107,177}

\begin{tikzpicture}[font=\sffamily\fontsize{7}{8}\selectfont]
  \def\OPC{.60}
\begin{groupplot}[
  group style={group size=2 by 2, horizontal sep=0.62in, vertical sep=0.75in},
  width=1.65in,
  height=1.30in,
  scale only axis=true,
  axis line style={draw=black!75, line width=0.45pt},
  tick style={draw=black!65, line width=0.35pt},
  tick align=outside,
  major tick length=2pt,
  %tick label style={font=\scriptsize},
  %label style={font=\scriptsize},
  %title style={%font=\footnotesize,
 %   yshift=1pt},
  grid=major,
  major grid style={draw=black!10, line width=0.30pt},
  minor tick num=0,    major tick length=0pt,
  scaled ticks=false,
  enlargelimits=false,
  %title style={yshift=-.1in},
]

% ------------------------------------------------------------------
% (a) MAP-alignment--novelty frontier.
% Labels are attached to the plotted points by nodes near coords;
% no axis coordinates are repeated in separate \node commands.
% ------------------------------------------------------------------
\nextgroupplot[
  title={(a) Excess alignment--novelty frontier},
  xlabel={Mean row novelty},
  ylabel={Excess MAP-alignment fraction $\Lambda$},
  xmin=-0.025, xmax=0.67,
  ymin=-0.08, ymax=1.20,
  xtick={0,0.2,0.4,0.6},
  ytick={0,0.25,0.50,0.75,1.00},
]
% Independence null: Lambda = 0.
%\addplot[black!55, dashed, thin, no marks, forget plot]
 % coordinates {(-0.025,0) (0.67,0)};

% Original-row reference: Lambda = 1.
%\addplot[black!55, solid, thin, no marks, forget plot]
%  coordinates {(-0.025,1) (0.67,1)};
\addplot[
  only marks, mark=*, mark size=1.85pt, lsmcol,
  error bars/.cd, x dir=both, y dir=both, x explicit, y explicit,
  error bar style={opacity=\OPC,line width=0.55pt},
  error mark options={rotate=90, mark size=1.5pt, line width=0.55pt},
  /pgfplots/.cd,
  nodes near coords={LSM},
  nodes near coords style={font=\scriptsize, text=lsmcol, anchor=east, xshift=-2pt, yshift=1pt}
] table[
  x=mean_row_novelty, y=mean_excess_alignment,
  x error plus expr=\thisrow{novelty_bootstrap_ci_high}-\thisrow{mean_row_novelty},
  x error minus expr=\thisrow{mean_row_novelty}-\thisrow{novelty_bootstrap_ci_low},
  y error plus expr=\thisrow{excess_alignment_ci_high}-\thisrow{mean_excess_alignment},
  y error minus expr=\thisrow{mean_excess_alignment}-\thisrow{excess_alignment_ci_low},
  col sep=comma
] {\SumLSM};

\addplot[
  only marks, mark=square*, mark size=1.70pt, chowcol,
  error bars/.cd, x dir=both, y dir=both, x explicit, y explicit,
  error bar style={line width=0.55pt},
  error mark options={rotate=90, mark size=1.5pt, line width=0.55pt},
  /pgfplots/.cd,
  nodes near coords={Chow--Liu},
  nodes near coords style={font=\scriptsize, text=chowcol, anchor=west, xshift=2pt, yshift=-2pt}
] table[
  x=mean_row_novelty, y=mean_excess_alignment,
  x error plus expr=\thisrow{novelty_bootstrap_ci_high}-\thisrow{mean_row_novelty},
  x error minus expr=\thisrow{mean_row_novelty}-\thisrow{novelty_bootstrap_ci_low},
  y error plus expr=\thisrow{excess_alignment_ci_high}-\thisrow{mean_excess_alignment},
  y error minus expr=\thisrow{mean_excess_alignment}-\thisrow{excess_alignment_ci_low},
  col sep=comma
] {\SumChow};

\addplot[
  only marks, mark=triangle*, mark size=1.85pt, ctgancol,
  error bars/.cd, x dir=both, y dir=both, x explicit, y explicit,
  error bar style={line width=0.55pt},
  error mark options={rotate=90, mark size=1.5pt, line width=0.55pt},
  /pgfplots/.cd,
  nodes near coords={CTGAN},
  nodes near coords style={font=\scriptsize, text=ctgancol, anchor=west, xshift=2pt, yshift=2pt}
] table[
  x=mean_row_novelty, y=mean_excess_alignment,
  x error plus expr=\thisrow{novelty_bootstrap_ci_high}-\thisrow{mean_row_novelty},
  x error minus expr=\thisrow{mean_row_novelty}-\thisrow{novelty_bootstrap_ci_low},
  y error plus expr=\thisrow{excess_alignment_ci_high}-\thisrow{mean_excess_alignment},
  y error minus expr=\thisrow{mean_excess_alignment}-\thisrow{excess_alignment_ci_low},
  col sep=comma
] {\SumCTGAN};

\addplot[
  only marks, mark=diamond*, mark size=1.75pt, basecol,
  error bars/.cd, x dir=both, y dir=both, x explicit, y explicit,
  error bar style={line width=0.55pt},
  error mark options={rotate=90, mark size=1.5pt, line width=0.55pt},
  /pgfplots/.cd,
  nodes near coords={Independent},
  nodes near coords style={font=\scriptsize, text=basecol, anchor=east, xshift=-2pt, yshift=-2pt}
] table[
  x=mean_row_novelty, y=mean_excess_alignment,
  x error plus expr=\thisrow{novelty_bootstrap_ci_high}-\thisrow{mean_row_novelty},
  x error minus expr=\thisrow{mean_row_novelty}-\thisrow{novelty_bootstrap_ci_low},
  y error plus expr=\thisrow{excess_alignment_ci_high}-\thisrow{mean_excess_alignment},
  y error minus expr=\thisrow{mean_excess_alignment}-\thisrow{excess_alignment_ci_low},
  col sep=comma
] {\SumBase};

\addplot[
  only marks, mark=*, mark size=1.65pt, ctrlcol,
  error bars/.cd, x dir=both, y dir=both, x explicit, y explicit,
  error bar style={line width=0.50pt},
  error mark options={rotate=90, mark size=1.4pt, line width=0.50pt},
  /pgfplots/.cd,
  nodes near coords={Original (control)},
  nodes near coords style={font=\scriptsize, text=ctrlcol, anchor=west, xshift=2pt, yshift=2pt}
] table[
  x=mean_row_novelty, y=mean_excess_alignment,
  x error plus expr=\thisrow{novelty_bootstrap_ci_high}-\thisrow{mean_row_novelty},
  x error minus expr=\thisrow{mean_row_novelty}-\thisrow{novelty_bootstrap_ci_low},
  y error plus expr=\thisrow{excess_alignment_ci_high}-\thisrow{mean_excess_alignment},
  y error minus expr=\thisrow{mean_excess_alignment}-\thisrow{excess_alignment_ci_low},
  col sep=comma
] {\SumCtrl};

% ------------------------------------------------------------------
% (b) Signed MAP-alignment gap. No mean triangles are drawn.
% Boxes show Q1--Q3, black line is median, whiskers use 1.5 IQR,
% and open circles are points beyond the whiskers.
% ------------------------------------------------------------------
\nextgroupplot[
  title={(b) Signed MAP-alignment gap},
  ylabel={$\Upsilon_{\mathrm{syn}}-\Upsilon_{\mathrm{control}}$},
  xmin=0.55, xmax=4.45,
  ymin=-0.075, ymax=0.045,
  xtick={1,2,3,4},
  xticklabels={LSM,Chow--Liu,CTGAN,Independent},
  x tick label style={rotate=25, anchor=east,yshift=-.05in},
  ytick={-0.06,-0.04,-0.02,0,0.02,0.04},
  yticklabel style={/pgf/number format/fixed,/pgf/number format/precision=2},
]
\ControlledBoxPlot{1}{lsmcol}{\WaveLSM}{signed_gap_vs_control}
\ControlledBoxPlot{2}{chowcol}{\WaveChow}{signed_gap_vs_control}
\ControlledBoxPlot{3}{ctgancol}{\WaveCTGAN}{signed_gap_vs_control}
\ControlledBoxPlot{4}{basecol}{\WaveBase}{signed_gap_vs_control}
\addplot[black!55, thin, no marks, forget plot] coordinates {(0.55,0) (4.45,0)};

% ------------------------------------------------------------------
% (c) Row novelty. Same explicit box construction.
% ------------------------------------------------------------------
\nextgroupplot[
  title={(c) Row novelty across waves},
  ylabel={Row novelty},
  xmin=0.55, xmax=5.45,
  ymin=-0.025, ymax=0.68,
  xtick={1,2,3,4,5},
  xticklabels={LSM,Chow--Liu,CTGAN,Independent,Original (control)},
  x tick label style={rotate=25, anchor=east,yshift=-.05in},
  ytick={0,0.2,0.4,0.6},
]
\ControlledBoxPlot{1}{lsmcol}{\WaveLSM}{row_novelty}
\ControlledBoxPlot{2}{chowcol}{\WaveChow}{row_novelty}
\ControlledBoxPlot{3}{ctgancol}{\WaveCTGAN}{row_novelty}
\ControlledBoxPlot{4}{basecol}{\WaveBase}{row_novelty}
\ControlledBoxPlot{5}{ctrlcol}{\WaveCtrl}{row_novelty}

% ------------------------------------------------------------------
% (d) Wave-wise retention.
% ------------------------------------------------------------------
\nextgroupplot[
  title={(d) Wave-wise MAP-alignment retention},
  xlabel={GSS wave [year]},
  ylabel={$\Upsilon_{\mathrm{synthetic}}/\Upsilon_{\mathrm{control}}$},
  xmin=1972, xmax=2024,
  ymin=0.905, ymax=1.055,
  xtick={1972,1980,1990,2000,2010,2020},
  xticklabel style={/pgf/number format/1000 sep=},
  ytick={0.92,0.96,1.00,1.04},
  legend style={font=\scriptsize, at={(0.02,1.01)}, anchor=north west,
    draw=none, fill=none, legend columns=2,
    /tikz/every even column/.append style={column sep=4pt}},
]
\addplot[lsmcol, opacity=\OPC,smooth, mark=*, mark size=0.95pt, line width=0.65pt]
  table[x=year, y=fidelity_retention, col sep=comma] {\WaveLSM};
\addlegendentry{LSM}
\addplot[chowcol, opacity=\OPC,smooth,mark=square*, mark size=0.90pt, line width=0.65pt]
  table[x=year, y=fidelity_retention, col sep=comma] {\WaveChow};
\addlegendentry{Chow--Liu}
\addplot[ctgancol, opacity=\OPC,smooth,mark=triangle*, mark size=1.00pt, line width=0.65pt]
  table[x=year, y=fidelity_retention, col sep=comma] {\WaveCTGAN};
\addlegendentry{CTGAN}
\addplot[basecol, opacity=\OPC,smooth,mark=diamond*, mark size=0.95pt, line width=0.65pt]
  table[x=year, y=fidelity_retention, col sep=comma] {\WaveBase};
\addlegendentry{Ind. baseline}
\addplot[black!55, thin, no marks, forget plot] coordinates {(1972,1) (2024,1)};

\end{groupplot}
\end{tikzpicture}

\caption{Four complementary views of the GSS benchmark: (a) the excess
MAP-alignment fraction
$\Lambda_g=(\bar{\Upsilon}_g-\bar{\Upsilon}_{\mathrm{ind}})/
(\bar{\Upsilon}_{\mathrm{control}}-\bar{\Upsilon}_{\mathrm{ind}})$
against mean row novelty, with nonparametric 95\% wave-bootstrap intervals; the
independent-marginal baseline is anchored at zero and the original (control)
at one; (b) wave-level signed MAP-alignment differences from the original (control);
(c) wave-level row novelty, defined as one minus mean nearest-real
similarity; and (d) MAP-alignment retention by GSS wave. The label
``Chow--Liu'' denotes the restricted Chow--Liu hybrid.}
\label{fig:gss_benchmark_four_panel}
\end{figure}


\section{Discussion}

Synthetic-data quality has two distinct axes. The original (control) preserves
MAP alignment but provides no novelty, whereas independent sampling creates new
rows while degrading conditional alignment. The desired regime combines
reference-level MAP alignment with separation from observed respondents.

The GSS benchmark spans surveys with different dimensionality and conditional
structure. Raw retention ratios are useful for within-wave comparison but are
not zero-anchored; the excess-alignment scale therefore measures recovery above
the dependence-destroying independent-marginal null. On this scale, LSM is
consistent with the original (control), whereas the other archived
configurations recover smaller fractions of the control-to-null increment. The
nearest-row diagnostic excludes simple resampling as the explanation. For
computational social science, the framework provides a model-agnostic validation
layer for synthetic survey populations used in sharing, replication, and
simulation.

\section{Limitations}

The identification results require strict positivity and compatibility of the
full conditional system. Separately fitted conditional learners need not be
exactly compatible with a single joint distribution. MAP alignment remains a
well-defined evaluator statistic in that setting, but interpreting $d_\mu$ as a
metric between underlying generative processes requires the stated compatibility
assumption.

The formal development is for finite product spaces. The GSS benchmark uses
categorical or discretely encoded variables, while the Gaussian examples use a
continuous-density analogue based on an essential supremum. A general treatment
of continuous and mixed-type variables, including the effects of discretization,
is not established here and is forthcoming.

The benchmark treats GSS records as unweighted observations and does not model
survey weights, strata, or primary sampling units. The reported quantities
therefore describe conditional structure in the analyzed records rather than
population-representative social-science estimands. Missing and nonresponse
entries were treated as missing rather than as substantive response categories;
the analysis did not statistically impute them. Results may nevertheless depend
on the prevalence and pattern of missingness across waves.

Computational cost grows with both the number of coordinates and the number of
evaluated records: the evaluator fits one conditional model per variable and
scores each record-coordinate pair. Parallelization made the 197--1,273-variable
GSS waves tractable, but repeated fitting, extensive tuning, and substantially
larger schemas may be costly.

\section{Future Work}

Immediate extensions are to formalize the continuous and mixed-type theory and
to incorporate survey weights and complex-design information into the reference
measure and aggregation scheme. Evaluator sensitivity should be assessed through
repeated train--test splits, repeated conditional-model fitting, and cross-fitted
variants when a separate evaluator-training set is unavailable. Empirical reports
should add exact duplicate rates and upper-tail nearest-row similarities, rather
than relying only on the mean novelty statistic. Future analyses should also test
sensitivity to missingness patterns and variable discretization, together with
computational approximations for very high-dimensional data.

\section{Data and Code Availability}

The GSS data are publicly available from NORC~\cite{gss2024}. The LSYNTH
implementation and reproducible examples are available in the public repository
at \url{https://github.com/zeroknowledgediscovery/lsynth}. The Python package is
distributed through PyPI and can be installed with \texttt{pip install lsynth}.

\bibliographystyle{splncs04}
\bibliography{lsyn}

\end{document}
